%!TeX program = xelatex
\documentclass[12pt,hyperref,a4paper,UTF8]{ctexart}
\usepackage{SEUReport}

%%-------------------------------正文开始---------------------------%%
\begin{document}

%%-----------------------封面--------------------%%
\cover

%%------------------摘要-------------%%
%\begin{abstract}
%
%在此填写摘要内容
%
%\end{abstract}

\thispagestyle{empty} % 首页不显示页码

%%--------------------------目录页------------------------%%
\newpage
\tableofcontents

%%------------------------正文页从这里开始-------------------%
\newpage

%%可选择这里也放一个标题
%\begin{center}
%    \title{ \Huge \textbf{{标题}}}
%\end{center}

\section{机器人总体概述}
\subsection{机器人设计目标与应用场景}
本机器人面向家庭服务场景进行设计,目标是在复杂、动态且非结构化的室内环境中,实现自主移动、环境感知、人机交互与目标操作的一体化协同。围绕RoboCup@Home比赛对机器人综合能力的要求,系统设计强调感知、决策、控制与执行的紧密结合,使机器人能够在家庭、实验室模拟家庭环境等典型室内场景中完成迎宾接待、语音交互、目标搜索、物品递送、分类搬运及基础服务操作等任务,从而具备较强的环境适应能力与任务执行能力。
\subsection{面向 RoboCup@Home赛项的任务定位}
面向RoboCup@Home赛项,本机器人定位为一套具备自主导航、自然交互、目标识别与抓取操作能力的综合型家庭服务机器人平台。结合比赛中Human Robot Interaction、Pick and Place、Doing Laundry和Restaurant等典型任务要求,机器人需要完成“感知环境—理解指令—自主移动—执行操作—反馈结果”的完整任务闭环,因此在系统设计上重点突出移动底盘、视觉感知、语音交互、机械臂与机械手等功能模块的协同集成,以满足家庭服务场景下多类型任务的执行需求。

\subsection{机器人整体外观}
下面是我们队伍设计的机器人整体外观示意图,
展示了底盘、躯干、传感器、机械臂与机械手的布局与结构设计。

\section{机器人本体结构设计}
\subsection{机械结构组成与整体布局}
机器人整体采用移动操作一体化的机械结构设计,主要由移动底盘、躯干支撑结构、传感器模块、机械臂及机械手等部分组成。底盘作为整机的运动与承载平台,承担机器人移动、供电以及上层结构支撑等功能；躯干部分采用铝方管搭建,作为整机主体框架,用于连接底盘与各功能模块,同时兼顾结构强度、安装灵活性与后期扩展性。机械臂通过铝方管支架安装于底盘上方,末端配备二指夹爪,用于完成比赛中的抓取、搬运与放置等操作任务。

在整体布局上,机器人采用分层集成的结构方案。底盘前方布置三维激光雷达,用于环境感知、障碍物检测与导航定位；机械臂安装于整机上部,以保证较好的作业高度与操作范围。为提升视觉感知能力,机械臂本体配备臂上相机,用于目标识别与抓取过程中的近距离视觉感知；机械臂下方额外安装一台相机,用于弥补上方相机在近距离观察和俯视场景中的视野盲区；此外,底盘后方还通过支架安装一台相机,用于补充机器人后方及整体环境的视觉感知范围。通过上述布局,机器人形成了以底盘运动层、躯干支撑层、操作执行层和多视角感知层相结合的整体结构,能够较好满足家庭服务机器人在导航、感知与操作等方面的综合需求。
\subsection{尺寸、重量、电源等基本参数}
机器人整机长\_\_\_mm、宽\_\_\_mm、高\_\_\_mm,整机重量约为\_\_\_kg。电源系统采用\_\_\_电池供电,额定电压为\_\_\_V,容量为\_\_\_Ah,可满足机器人底盘驱动、机械臂执行机构、传感器及主控计算平台的稳定运行,续航时间约为\_\_\_小时。整体参数设计兼顾结构紧凑性、运行稳定性与比赛场景适应性,能够满足家庭服务机器人在室内环境中的移动、感知与操作需求。
\subsection{急停按钮、启动按钮与安全设计}

\subsection{学校及队伍标识设计}

\section{移动底盘硬件与驱动系统}
\subsection{底盘结构与运动方式}
本机器人采用松灵 TRACER V2 作为移动底盘。该底盘采用两轮差速驱动结构,具有运动灵活、控制方式成熟、系统集成方便等特点,能够满足家庭服务机器人在室内环境中的运动需求。通过左右驱动轮速度差,机器人可实现前进、后退、转向及原地旋转等基本运动,在门口、桌边、家具间通道等相对狭窄的场景中具有较好的机动性与通过性。同时,差速驱动方式便于与自主导航、路径跟踪和局部避障等上层算法进行结合,有利于机器人在比赛任务中完成目标点停靠、位置调整与近距离操作前的姿态对准。

\subsection{底盘驱动系统}
本机器人底盘采用松灵 TRACER V2 平台,并直接使用其原生驱动与控制系统。底盘与主机之间通过 CAN 总线进行通信,由主机运行 ROS 节点完成运动指令下发、底盘状态读取以及与上层导航控制模块的数据交互。该方式能够较好实现底盘驱动层与机器人整体软件系统的集成,使底盘运动控制能够与定位、路径规划、局部避障等功能协同工作。

在实际系统中,主机侧 ROS 节点根据上层任务规划与导航控制结果生成线速度和角速度指令,并通过 CAN 通信发送到底盘控制系统,由底盘执行相应的运动控制；同时,底盘运行状态与运动反馈信息也可回传至主机,为上层控制提供必要的状态支撑。该驱动方式具有接口清晰、实时性较好、系统集成方便等优点,能够满足比赛环境下机器人对运动稳定性、控制可靠性以及软件扩展性的要求。

\section{机械臂与机械手配置}
\subsection{机械臂型号、自由度与安装方式}

本机器人采用松灵 PiPER 机械臂作为操作执行机构。该机械臂为六自由度串联结构,具有较好的空间可达性与操作灵活性,能够满足家庭服务机器人在抓取、搬运与放置等任务中的使用需求。机械臂安装于机器人上部铝方管支撑结构上,通过底座与支架进行刚性连接,从而保证移动过程中机械臂工作的稳定性,并为末端在桌面及柜体附近区域的操作提供合适的作业空间。

\subsection{机械手结构、适用物体类型与抓取方式}

机械臂末端采用两指夹爪作为机械手,具有结构简单、控制直接和可靠性较高等特点,适合比赛环境下对多类日常物品进行抓取。结合赛项任务需求,该机械手主要面向餐具、瓶罐、盒装物品、小型生活用品以及衣物等目标开展操作。对于刚性较好的物体,系统主要采用直接夹持的方式完成抓取与搬运；对于柔软衣物,则先将夹爪张开至合适宽度并贴近衣物表面,在夹爪闭合过程中利用摩擦力带动衣物向内形成褶皱,从而获得较稳定的抓取点。该方案已经过实际测试验证,能够较好提升衣物抓取的成功率与稳定性。

\subsection{机械臂控制器与执行驱动方案}

PiPER 机械臂采用集成式控制器设计,通过 CAN 线经 USB 转接模块与主机连接,由主机侧程序完成机械臂控制与状态读取。在执行方案上,系统主要使用 MoveIt 进行机械臂运动规划与轨迹执行,以完成抓取、放置等任务过程中的位姿求解和关节运动控制；同时结合官方 SDK 获取关节状态及关节力信息,为抓取接触检测和执行过程监测提供支持。该控制方式便于机械臂与视觉感知、任务规划及整机系统进行协同集成,能够满足比赛任务中的操作需求。

\section{感知硬件配置}

\subsection{主控设备配置}

本机器人主控设备采用队员使用的笔记本电脑,作为整机感知、控制与任务执行的核心计算平台。主机主要负责运行机器人系统中的 ROS 节点、导航与感知算法、机械臂控制程序以及人机交互相关功能模块,从而实现各硬件设备之间的数据处理与协同控制。与此同时,系统外接显示屏用于交互信息的显示,可在比赛过程中向操作对象展示提示信息、任务状态或识别结果,从而增强机器人在人机交互场景中的信息表达能力。

\subsection{激光雷达配置}

机器人配备 Livox Mid-360 激光雷达作为主要环境感知传感器。该雷达主要用于获取周围环境的三维点云信息,是机器人完成环境建图、定位与自主导航的重要感知硬件。在系统中,Mid-360 所提供的点云数据可用于支撑 SLAM 算法运行,帮助机器人实现对比赛场地环境的感知与理解,并为后续路径规划、障碍物检测和运动控制提供基础信息支持。

\subsection{相机配置}

本机器人共配置三台相机,以满足不同任务场景下的视觉感知需求。其中,机械臂上的相机采用 RealSense L515,主要用于抓取任务中的物品感知,为目标识别、定位及抓取过程提供近距离视觉信息；机器人顶部相机采用 RealSense D435,主要服务于人机交互任务,用于进行人物相关特征的识别,如人脸、人体以及手势等；机器人正面中部同样配置一台 RealSense D435,相较于顶部相机更关注环境中的低矮目标,主要用于垃圾桶、洗衣机等目标的识别与观测。通过多位置相机协同配置,机器人能够兼顾人机交互、环境理解与操作任务中的视觉感知需求。

\subsection{麦克风、扬声器}

在语音交互硬件方面,机器人采用外接收音麦克风作为语音输入设备,用于采集用户语音指令和交互信息；扬声器则采用笔记本电脑自带扬声器,用于语音播报、任务反馈及提示信息输出。该配置能够满足比赛过程中基本的人机语音交互需求,并与机器人视觉感知及任务执行模块共同构成完整的交互链路。


\section{物品抓取软件算法}

针对比赛中常见的餐具、瓶罐、盒装物品及其他日常用品抓取任务,本机器人构建了由目标识别、三维感知、抓姿生成与运动执行组成的物品抓取软件算法流程。考虑到比赛环境中待抓取目标既包含类别相对固定的常见物品,也可能存在外观差异较大、语义描述更加灵活的目标对象,系统在目标识别阶段采用闭集识别与开放语义识别相结合的方式。其中,YOLOv11 主要用于完成常见目标物品的快速检测,获得目标在图像中的位置与类别信息；开放语义物品识别模型则用于对目标语义进行补充匹配与进一步确认,从而提升系统对复杂目标和描述性目标的识别能力。通过上述方法,机器人能够在较复杂的场景中更准确地确定待抓取物体,为后续抓取规划提供可靠输入。

在完成目标识别后,系统进一步结合深度信息提取目标物品附近的局部点云,并以此作为抓取姿态生成模块的输入。针对局部点云,系统采用 GraspNet 进行抓姿推理,生成候选抓取位姿,再结合机械臂当前状态与执行可达性筛选出合适的抓取姿态。随后,系统利用 MoveIt 对机械臂从当前位姿到目标抓取位姿的运动过程进行轨迹规划,并控制机械臂逐步接近目标位置,在到达预定抓取位姿后闭合夹爪完成物品抓取。整体来看,该方法将二维目标识别、三维点云理解、数据驱动抓姿生成与机械臂运动规划有机结合,形成了较为完整的自主抓取技术链路,能够较好适应比赛任务中目标种类多样、摆放状态复杂及抓取过程要求连续稳定的应用需求。
\section{地图构建、定位与自主导航}
\subsection{场景地图构建方法}

\subsection{定位与重定位方法}

\subsection{全局路径规划与局部避障策略}

\section{视觉感知与图像识别技术}
\subsection{人体检测与人物识别}

\subsection{物体检测、识别与分类}

\subsection{手势、挥手、指向等行为识别}

\subsection{场景理解与目标定位}

\subsection{面向抓取任务的视觉伺服或位姿估计方法}

\subsection{不同光照和遮挡条件下的鲁棒性设计}

\section{语音交互与人机交互系统}
\subsection{语音采集与唤醒方式}

\subsection{语音识别与语义理解方法}

\subsection{对话管理与任务指令解析}

%%----------- 参考文献 -------------------%%
%在reference.bib文件中填写参考文献,此处自动生成

\reference


\end{document}
